\clearpage
\section{Design}

The design preserves the \textbf{domain structure} of assignment 3 and adds only
the distributed mechanisms that are necessary for clustering. The main goal is
to keep the \textbf{alarm semantics unchanged} while making the runtime satisfy
three distributed constraints: \textbf{multiple nodes}, \textbf{message-only
communication}, and \textbf{safe recovery after control-unit recreation}.

\subsection{Role-based node architecture}

Each \textbf{cluster node} is started with exactly one role:
\begin{itemize}
    \item \texttt{control-unit}
    \item \texttt{keypad}
    \item \texttt{sensor}
\end{itemize}

This is more explicit than a monolithic startup. The selected role determines
the node's main responsibility and which role-specific actors are started locally.
The chosen deployment rule is therefore \textbf{one node, one main role}. This
is not a Pekko limitation, but a deliberate design choice that makes the
distributed architecture easier to verify, demonstrate, and explain.

\subsection{Single logical control unit}

The central design decision is that the \textbf{alarm state machine} still has
one logical owner: \texttt{ControlUnitActor}. In the clustered version it is not
spawned as an ordinary local actor, but exposed through a \textbf{cluster
singleton}. This guarantees that, even if multiple control-unit nodes are
started, only \textbf{one active control entity} manages the alarm state at a time.

The solution adopted in the final project is therefore to keep the
\textbf{Control Unit} as the only authority, expose it as a \textbf{cluster
singleton}, and force it to restart in \texttt{STARTUP\_RECOVERY} after
recreation.

This choice is aligned with the domain model. The system does not replicate or
merge alarm state across nodes, so allowing multiple active control units would
create ambiguity. The singleton makes the ownership of the alarm state
\textbf{explicit} and \textbf{cluster-wide}.

\subsection{Discovery of distributed actors}

In assignment 3 the actors could receive direct references during local
construction. In assignment 4 this is no longer sufficient, because keypad and
sensor nodes may run on different cluster nodes.

The adopted solution is \textbf{receptionist-based discovery}:
\begin{itemize}
    \item the control unit registers itself as a discoverable service;
    \item keypads and sensors subscribe to that service and keep the discovered reference updated;
    \item sirens register themselves, and the control unit subscribes to the available siren instances.
\end{itemize}

This solution keeps the system \textbf{loosely coupled}. Peripheral actors do not
need to know a concrete node address, and they do not need custom logic for
which control-unit node is currently hosting the singleton.

\subsection{Recovery semantics}

The main behavioral extension with respect to assignment 3 is the state
\texttt{STARTUP\_RECOVERY}. After creation or recreation, the control unit:
\begin{itemize}
    \item does \textbf{not} assume the system is armed;
    \item does \textbf{not} assume the system is disarmed;
    \item ignores or logs sensor activations;
    \item accepts only a \textbf{correct PIN} to return to \texttt{DISARMED}.
\end{itemize}

This matches the specification exactly. State is intentionally \textbf{not
persisted}, so the recovery policy cannot reconstruct the previous condition.
The correct design response is therefore \textbf{safe recovery}, not optimistic
recovery.

\subsection{Trade-offs}

This solution prefers \textbf{clarity} over maximum generality. It does not add
persistence, replicated state, or advanced failover policies, because those
features are outside the assignment goals. The chosen design is sufficient to
show:
\begin{itemize}
    \item distributed startup;
    \item remote discovery;
    \item single ownership of state;
    \item message-only interaction;
    \item safe recovery after recreation.
\end{itemize}

In other words, the design intentionally concentrates on the concepts that are
most relevant for this assignment: \textbf{where actors run}, \textbf{how they
find each other}, \textbf{who owns the alarm state}, and \textbf{how the system
behaves after failure}.
